\section{问题分析}

四问共享“附件数据重构—内部热湿传递—表面对流交换—数值求解—结果验证”的主链条，但各问新增困难不同。问题一解决基础的时变边界扩散；问题二处理系数随解变化引起的非线性；问题三引入全域阈值首达事件；问题四再处理随时间收缩的空间区域。按照这一依赖关系建立同一套守恒离散程序，可以避免各问采用不同数值口径造成的结果不可比。总体技术路线如\cref{fig:roadmap}所示。

\begin{figure}[!htbp]
  \centering
  \includegraphics[width=0.98\textwidth]{overall_roadmap.pdf}
  \caption{四问递进建模与验证技术路线}
  \label{fig:roadmap}
\end{figure}

\subsection{问题一的分析}

问题一在固定半径及主要热物性下，建立一维径向热湿传递模型，保留水分扩散系数 $D(C)$ 的非线性。由附件重构时变环境，在中心和表面分别设置对称与 Robin 边界，采用有限体积法求解预热阶段的温湿分布，为后续各问提供基础。


\subsection{问题二的分析}

问题二引入随温湿状态变化的密度、比热、导热及扩散系数，形成非线性热湿耦合模型。采用分步 Picard 迭代更新温度、水分与物性，获得同一时间层的自洽解；比较中心与表面的响应，分析升温和内部水分扩散对干燥过程的影响。


\subsection{问题三的分析}

问题三以全域最大水分浓度首次达到 $0.15~\mathrm{kg/kg}$ 为干燥完成判据，在问题二模型上延长计算时间。通过全过程剖面验证最大含水位置，确定控制达标的部位，并在阈值附近进行时间插值，以准确定位烘干结束时刻。


\subsection{问题四的分析}

问题四考虑药材均匀径向收缩，利用 $\xi=r/R(t)$ 将移动域映射至固定区间，求解变物性热湿迁移过程。另设同物性固定半径模型以分离收缩贡献，并采用物理坐标移动网格与 COMSOL 有限元计算交叉验证，检验移动边界求解的可靠性。
